\section{Implementation Details}
\label{sec:implementation}

WuYa is implemented as a modular multi-agent system. This section provides technical details on the implementation stack, data models, and key engineering decisions.

\subsection{Technical Stack}

The system is built using the following components:

\begin{itemize}
    \item \textbf{LLM Backend}: GPT-4 or Claude-3 for \subagent{} reasoning; configurable to use local models (Llama-3, Mistral) for privacy-sensitive deployments.
    \item \textbf{Vector Database}: Milvus for \rag{} indexing and retrieval; supports dense embeddings (OpenAI text-embedding-3) and sparse retrieval (BM25 hybrid).
    \item \textbf{Journal Database}: PostgreSQL for structured journal metadata; includes SpringRank tier information for 14,590 journals across all disciplines.
    \item \textbf{Citation Network}: Neo4j for paper citation graph; enables efficient frontier construction for \dea{} analysis.
    \item \textbf{Message Queue}: Redis for asynchronous task scheduling; enables parallel \subagent{} execution and result aggregation.
    \item \textbf{API Layer}: FastAPI for RESTful endpoints; supports webhook callbacks for async operations.
\end{itemize}

\subsection{Core Data Models}

The system operates on the following data structures:

\begin{verbatim}
Paper {
  id: UUID
  content: Text
  metadata: {
    title: String
    authors: [String]
    abstract: Text
    keywords: [String]
    discipline: String
    citations: [UUID]
    cited_by: [UUID]
  }
  evaluations: [EvaluationResult]
}

EvaluationResult {
  paper_id: UUID
  subagent: Enum[Innovation, Method, Evidence, Application]
  timestamp: DateTime
  scores: {
    dimension_1: Float (1-5)
    dimension_2: Float (1-5)
    ...
  }
  narrative: Text
  rag_citations: [Citation]
}

ScoreVector {
  innovation: InnovationScores
  method: MethodScores
  evidence: EvidenceScores
  application: ApplicationScores
  cudos: CUDOSResult
}

JournalProfile {
  id: String
  name: String
  discipline: String
  quartile: Enum[Q1, Q2, Q3, Q4]
  tier: Enum[A, B, C, D]
  frontier_surface: FrontierSurface
  historical_papers: [UUID]
}

FrontierSurface {
  journal_id: String
  nodes: [FrontierNode]
  updated_at: DateTime
  confidence: Float
}

FrontierNode {
  dimension: String
  preference: String
  evidence: [String]
  confidence: Float
}
\end{verbatim}

\subsection{Router Orchestration Logic}

The Router implements the two-phase workflow as follows:

\begin{algorithm}[H]
\caption{Router Orchestration}
\label{alg:router}
\begin{algorithmic}
\STATE Input: paper, user\_intent
\STATE parsed $\leftarrow$ parse\_paper(paper)
\STATE discipline $\leftarrow$ identify\_discipline(parsed)

\COMMENT{Phase 1: CUDOS Gatekeeping}
\STATE cudos\_result $\leftarrow$ await call\_subagent("cudos", parsed)
\IF{not cudos\_result.gate\_pass}
    \STATE RETURN generate\_veto\_report(cudos\_result)
\ENDIF

\COMMENT{Phase 2: Parallel Evaluation}
\STATE evaluation\_tasks $\leftarrow$ [
    call\_subagent("innovation", parsed),
    call\_subagent("method", parsed),
    call\_subagent("evidence", parsed),
    call\_subagent("application", parsed)
]
\STATE scores $\leftarrow$ await gather(evaluation\_tasks)
\STATE score\_vector $\leftarrow$ aggregate\_scores(scores)

\COMMENT{Paper Localization}
\STATE path\_a $\leftarrow$ llm\_mapper(score\_vector, discipline, get\_prior(discipline))
\IF{user\_intent.has\_target\_journal}
    \STATE path\_b $\leftarrow$ dea\_analysis(score\_vector, user\_intent.journal)
\ELSE
    \STATE path\_b $\leftarrow$ None
\ENDIF

\IF{path\_b and not consistent(path\_a, path\_b)}
    \STATE rag\_explanation $\leftarrow$ await generate\_rag\_explanation(path\_a, path\_b)
\ENDIF

\IF{user\_intent.type == "recommendation"}
    \STATE RETURN generate\_recommendation(score\_vector, path\_a, path\_b, rag\_explanation)
\ELSE
    \STATE RETURN generate\_review\_report(score\_vector, path\_a, path\_b, rag\_explanation)
\ENDIF
\end{algorithmic}
\end{algorithm}

\subsection{RAG Implementation}

The \rag{} layer indexes canonical texts organized by theorist and concept:

\begin{itemize}
    \item \textbf{Primary sources}: Full texts of Popper, Kuhn, Lakatos, Merton, Fisher, Campbell, Pearl, Bush, Mokyr.
    \item \textbf{Secondary sources}: Commentary and application of these theories in modern contexts.
    \item \textbf{Chunking strategy}: Semantic chunking at paragraph level (avg. 300 tokens); overlap of 50 tokens for context.
    \item \textbf{Retrieval}: Hybrid dense (embedding similarity) + sparse (BM25) retrieval; reranked by cross-encoder.
    \item \textbf{Trigger integration}: \subagents{} invoke RAG through a unified interface with query rewriting and result filtering.
\end{itemize}

\subsection{Engineering Considerations}

\textbf{Caching}: Results for identical (paper\_id, score\_vector) pairs are cached for 24 hours to reduce LLM calls. Frontier Surface Library is cached with TTL of 1 hour.

\textbf{Error Handling}: Each \subagent{} has a timeout of 60 seconds; on failure, the Router retries twice before returning partial results with an error flag.

\textbf{Observability}: All \subagent{} calls are logged with trace IDs for debugging. Key metrics include: \subagent{} latency, RAG hit rate, evaluation consistency across runs.

\textbf{Scalability}: The system supports horizontal scaling through stateless \subagent{} workers; Redis queue distributes tasks across workers. Target throughput: 100 papers/minute with 8 workers.
